% Les cinq images des acronymes (section non numérotée, à inclure par % Les cinq images des acronymes (section non numérotée, à inclure par % Les cinq images des acronymes (section non numérotée, à inclure par % Les cinq images des acronymes (section non numérotée, à inclure par \input{easter_eggs_figures}).
% Figures produites par code/easter_egg_images.py (matplotlib/numpy/scipy) ; chaque grandeur tracée est calculée
% à partir de sa formule. Les auto-contrôles numériques (contrôles positifs et négatifs) sont consignés dans
% results/textbook_formalisation_2026/easter_eggs_selfcheck.json ; un contrôle en échec interrompt le script.
\section*{Les cinq images}
\addcontentsline{toc}{section}{Les cinq images}

Chacun des cinq noms courts de la section~\ref{ch5:acronymes} reçoit ici une image. Le niveau du théorème
(\tierA{}) porte sur l'énoncé Lean ré-exporté. Les courbes, les couleurs et les valeurs affichées sont des calculs
numériques \tierX{} : elles illustrent un cas particulier de l'énoncé, sans le démontrer ni l'étendre.

\begin{figure}[htbp]
\centering
\includegraphics[width=\textwidth]{egg_ALIX.pdf}
\caption{\lean{ALIX}, \emph{A Lattice Is eXtremal}. Le théorème \tierA{} (conditionnel au théorème amont
\lean{universal_energy_minimum} d'openai/math, généré par IA et non relu par des pairs) dit qu'à toute densité
$\rho>0$ le réseau triangulaire minimise l'énergie de tout potentiel admissible, parmi toutes les configurations
localement finies de densité~$\rho$. L'image \tierX{} prend le noyau gaussien $e^{-\pi\alpha r^2}$, $\alpha=0{,}4$,
à densité~1, et ne compare que des \emph{réseaux}. (a)~et~(b) : cellules de Voronoï colorées par la contribution
$e^{-\pi\alpha|p|^2}$ de chaque site~$p$ à l'énergie du site central ; $E_\alpha=1{,}50173$ (triangulaire)
contre $1{,}50388$ (carré). (c)~: écart $E_\alpha(\tau)-E_\alpha(e^{i\pi/3})$ sur le domaine fondamental des formes
de réseau, calculé par sommes thêta ; le minimum de la grille (au point équivalent $-\tfrac12+i\tfrac{\sqrt3}2$, à un pas de grille près) et cinq minimisations de Nelder--Mead (à $10^{-7}$ près) retombent sur
$\tau=e^{i\pi/3}$. En médaillon : $E_{\text{carré}}/E_{\text{tri}}-1>0$ pour
$\alpha\in[0{,}15;6]$, c'est-à-dire pour toute une gamme de densités.}
\label{fig:egg:ALIX}
\end{figure}

\begin{figure}[htbp]
\centering
\includegraphics[width=\textwidth]{egg_KAL.pdf}
\caption{\lean{KAL}, \emph{Kolmogorov Absorption Limit}. Le théorème \tierA{} : pour une chaîne sur
$\{0,\dots,N\}$ à bords absorbants, si $h$ est harmonique aux états intérieurs, $h(0)=0$, $h(N)=1$, et si la
variance à un pas de $h$ est strictement positive, alors $P^n(i,N)\to h(i)$ et $P^n(i,0)\to 1-h(i)$. Images
\tierX{}, par puissances exactes de la matrice de transition. (a)~Wright--Fisher neutre, $N=24$ : $P^n(i,N)$
converge vers $h(i)=i/N$ (bande de droite) ; écart maximal à $n=4096$ : $6\cdot10^{-16}$. (b)~Moran avec
sélection, $N=24$, $r=1{,}25$ : $P^n(i,N)$ tend vers $(1-r^{-i})/(1-r^{-N})$ (pointillés), écart
$7\cdot10^{-15}$ à $n=2^{22}$. L'harmonicité des deux $h$ est vérifiée en rationnels exacts (pour Wright--Fisher, à $N=8$).}
\label{fig:egg:KAL}
\end{figure}

\begin{figure}[htbp]
\centering
\includegraphics[width=\textwidth]{egg_ILENE.pdf}
\caption{\lean{ILENE}, \emph{Isothermal Law: Every Node Equivalent}. Le théorème \tierA{} : sur un graphe
isotherme ($W$ bistochastique) fortement connexe, en naissance--mort à sélection uniforme, un mutant unique de
fitness $r\neq1$ se fixe avec la probabilité $(1-r^{-1})/(1-r^{-n})$, quel que soit son sommet de départ. Images
\tierX{}, $r=2$, probabilités obtenues en résolvant le système linéaire sur les $2^n$ configurations (en flottants pour $n=8$, en rationnels pour $n=3$), et
contrôlées par la limite de $P^t$. (a)~Graphe pondéré à $n=8$ sommets, $W=\frac12\Pi_1+\frac13\Pi_2+\frac16\Pi_3$
(trois permutations sans point fixe) : les huit sommets donnent la valeur de Moran $0{,}50196$ à $2\cdot10^{-15}$ près.
(b)~L'étoile à 7 feuilles n'est pas isotherme et sort des hypothèses : centre $0{,}167$, feuilles $0{,}700$.
(c)~Calcul en rationnels pour $n=3$ : le 3-cycle donne $\frac47$ partout ; l'étoile donne $\frac5{12}$ au centre et
$\frac23$ aux feuilles, de moyenne $\frac7{12}\neq\frac47$.}
\label{fig:egg:ILENE}
\end{figure}

\begin{figure}[htbp]
\centering
\includegraphics[width=\textwidth]{egg_ALEXIS.pdf}
\caption{\lean{ALEXIS}, \emph{Averaged Laplace EXponentials Inherit Signs}. Le théorème \tierA{} : un mélange
positif d'exponentielles $t\mapsto\int e^{-t\,k(x)}m(x)\,d\mu(x)$ est un potentiel admissible, de classe
$C^\infty$ sur $]0,+\infty[$, positif, avec $(-1)^n g^{(n)}\ge0$. L'image \tierX{} en montre une instance, le
potentiel de Rytova--Keldysh ($r_0=1$), écrit comme mélange de gaussiennes
$V(r)=\int_0^\infty \frac{w(a)}{2a}\,e^{-r^2/4a}\,da$ de poids $w(a)=\frac1{r_0}e^{a/r_0^2}\operatorname{erfc}
(\sqrt a/r_0)>0$. (a)~Les contributions de douze tranches en $a$, empilées, redonnent la forme de Struve
$\frac{\pi}{2r_0}[\mathbf H_0-Y_0](r/r_0)$ (points) à moins de $3\cdot10^{-13}$ près en relatif ; l'identification avec la
forme de Struve est une preuve humaine, non vérifiée par machine \tierC{}. (b)~Le poids $w$, coloré par tranche.
(c)~$(-1)^n g^{(n)}(t)$, $g(t)=V(\sqrt t)$, pour $n=0,\dots,6$ : toutes les courbes sont positives, donc les
dérivées $g^{(n)}$ alternent de signe ; $n=1$ est recoupé par différences finies de la forme de Struve.}
\label{fig:egg:ALEXIS}
\end{figure}

\begin{figure}[htbp]
\centering
\includegraphics[width=\textwidth]{egg_AURELIA.pdf}
\caption{\lean{AURELIA}, \emph{Average Under Replication: Evolution Lifts It, Always}. Le théorème \tierA{} :
en sélection haploïde à fitness constantes, sans biais de transmission, la fitness moyenne ne décroît pas,
$\bar w\le\bar w'$ ; c'est un cas particulier du théorème fondamental de Fisher, déduit de l'équation de Price.
Images \tierX{}. (a)~Trois types, $w=(1;\,1{,}45;\,1{,}9)$ : sur le simplexe coloré par $\bar w$, six trajectoires
de $x_i'=x_iw_i/\bar w$ montent vers le meilleur type. (b)~Cas diploïde sous Hardy--Weinberg, avec surdominance
($w_{11}=0{,}78$, $w_{12}=1$, $w_{22}=0{,}62$) : les deux trajectoires montent sur $\bar w(p)$ jusqu'à
$p^*=0{,}6333$. Ce panneau relève du théorème diploïde \lean{wbar_increment} du chapitre~4 (H12) \tierA{}, pas de
\lean{AURELIA}. (c)~$\bar w$ au fil des générations, jamais décroissante. L'identité
$\bar w'-\bar w=\mathrm{Var}(w)/\bar w$ est vérifiée en rationnels exacts sur 400 cas aléatoires ; contrôle
négatif : avec un biais de mutation du meilleur type vers le moins bon, $\bar w$ baisse de $0{,}35$.}
\label{fig:egg:AURELIA}
\end{figure}
).
% Figures produites par code/easter_egg_images.py (matplotlib/numpy/scipy) ; chaque grandeur tracée est calculée
% à partir de sa formule. Les auto-contrôles numériques (contrôles positifs et négatifs) sont consignés dans
% results/textbook_formalisation_2026/easter_eggs_selfcheck.json ; un contrôle en échec interrompt le script.
\section*{Les cinq images}
\addcontentsline{toc}{section}{Les cinq images}

Chacun des cinq noms courts de la section~\ref{ch5:acronymes} reçoit ici une image. Le niveau du théorème
(\tierA{}) porte sur l'énoncé Lean ré-exporté. Les courbes, les couleurs et les valeurs affichées sont des calculs
numériques \tierX{} : elles illustrent un cas particulier de l'énoncé, sans le démontrer ni l'étendre.

\begin{figure}[htbp]
\centering
\includegraphics[width=\textwidth]{egg_ALIX.pdf}
\caption{\lean{ALIX}, \emph{A Lattice Is eXtremal}. Le théorème \tierA{} (conditionnel au théorème amont
\lean{universal_energy_minimum} d'openai/math, généré par IA et non relu par des pairs) dit qu'à toute densité
$\rho>0$ le réseau triangulaire minimise l'énergie de tout potentiel admissible, parmi toutes les configurations
localement finies de densité~$\rho$. L'image \tierX{} prend le noyau gaussien $e^{-\pi\alpha r^2}$, $\alpha=0{,}4$,
à densité~1, et ne compare que des \emph{réseaux}. (a)~et~(b) : cellules de Voronoï colorées par la contribution
$e^{-\pi\alpha|p|^2}$ de chaque site~$p$ à l'énergie du site central ; $E_\alpha=1{,}50173$ (triangulaire)
contre $1{,}50388$ (carré). (c)~: écart $E_\alpha(\tau)-E_\alpha(e^{i\pi/3})$ sur le domaine fondamental des formes
de réseau, calculé par sommes thêta ; le minimum de la grille (au point équivalent $-\tfrac12+i\tfrac{\sqrt3}2$, à un pas de grille près) et cinq minimisations de Nelder--Mead (à $10^{-7}$ près) retombent sur
$\tau=e^{i\pi/3}$. En médaillon : $E_{\text{carré}}/E_{\text{tri}}-1>0$ pour
$\alpha\in[0{,}15;6]$, c'est-à-dire pour toute une gamme de densités.}
\label{fig:egg:ALIX}
\end{figure}

\begin{figure}[htbp]
\centering
\includegraphics[width=\textwidth]{egg_KAL.pdf}
\caption{\lean{KAL}, \emph{Kolmogorov Absorption Limit}. Le théorème \tierA{} : pour une chaîne sur
$\{0,\dots,N\}$ à bords absorbants, si $h$ est harmonique aux états intérieurs, $h(0)=0$, $h(N)=1$, et si la
variance à un pas de $h$ est strictement positive, alors $P^n(i,N)\to h(i)$ et $P^n(i,0)\to 1-h(i)$. Images
\tierX{}, par puissances exactes de la matrice de transition. (a)~Wright--Fisher neutre, $N=24$ : $P^n(i,N)$
converge vers $h(i)=i/N$ (bande de droite) ; écart maximal à $n=4096$ : $6\cdot10^{-16}$. (b)~Moran avec
sélection, $N=24$, $r=1{,}25$ : $P^n(i,N)$ tend vers $(1-r^{-i})/(1-r^{-N})$ (pointillés), écart
$7\cdot10^{-15}$ à $n=2^{22}$. L'harmonicité des deux $h$ est vérifiée en rationnels exacts (pour Wright--Fisher, à $N=8$).}
\label{fig:egg:KAL}
\end{figure}

\begin{figure}[htbp]
\centering
\includegraphics[width=\textwidth]{egg_ILENE.pdf}
\caption{\lean{ILENE}, \emph{Isothermal Law: Every Node Equivalent}. Le théorème \tierA{} : sur un graphe
isotherme ($W$ bistochastique) fortement connexe, en naissance--mort à sélection uniforme, un mutant unique de
fitness $r\neq1$ se fixe avec la probabilité $(1-r^{-1})/(1-r^{-n})$, quel que soit son sommet de départ. Images
\tierX{}, $r=2$, probabilités obtenues en résolvant le système linéaire sur les $2^n$ configurations (en flottants pour $n=8$, en rationnels pour $n=3$), et
contrôlées par la limite de $P^t$. (a)~Graphe pondéré à $n=8$ sommets, $W=\frac12\Pi_1+\frac13\Pi_2+\frac16\Pi_3$
(trois permutations sans point fixe) : les huit sommets donnent la valeur de Moran $0{,}50196$ à $2\cdot10^{-15}$ près.
(b)~L'étoile à 7 feuilles n'est pas isotherme et sort des hypothèses : centre $0{,}167$, feuilles $0{,}700$.
(c)~Calcul en rationnels pour $n=3$ : le 3-cycle donne $\frac47$ partout ; l'étoile donne $\frac5{12}$ au centre et
$\frac23$ aux feuilles, de moyenne $\frac7{12}\neq\frac47$.}
\label{fig:egg:ILENE}
\end{figure}

\begin{figure}[htbp]
\centering
\includegraphics[width=\textwidth]{egg_ALEXIS.pdf}
\caption{\lean{ALEXIS}, \emph{Averaged Laplace EXponentials Inherit Signs}. Le théorème \tierA{} : un mélange
positif d'exponentielles $t\mapsto\int e^{-t\,k(x)}m(x)\,d\mu(x)$ est un potentiel admissible, de classe
$C^\infty$ sur $]0,+\infty[$, positif, avec $(-1)^n g^{(n)}\ge0$. L'image \tierX{} en montre une instance, le
potentiel de Rytova--Keldysh ($r_0=1$), écrit comme mélange de gaussiennes
$V(r)=\int_0^\infty \frac{w(a)}{2a}\,e^{-r^2/4a}\,da$ de poids $w(a)=\frac1{r_0}e^{a/r_0^2}\operatorname{erfc}
(\sqrt a/r_0)>0$. (a)~Les contributions de douze tranches en $a$, empilées, redonnent la forme de Struve
$\frac{\pi}{2r_0}[\mathbf H_0-Y_0](r/r_0)$ (points) à moins de $3\cdot10^{-13}$ près en relatif ; l'identification avec la
forme de Struve est une preuve humaine, non vérifiée par machine \tierC{}. (b)~Le poids $w$, coloré par tranche.
(c)~$(-1)^n g^{(n)}(t)$, $g(t)=V(\sqrt t)$, pour $n=0,\dots,6$ : toutes les courbes sont positives, donc les
dérivées $g^{(n)}$ alternent de signe ; $n=1$ est recoupé par différences finies de la forme de Struve.}
\label{fig:egg:ALEXIS}
\end{figure}

\begin{figure}[htbp]
\centering
\includegraphics[width=\textwidth]{egg_AURELIA.pdf}
\caption{\lean{AURELIA}, \emph{Average Under Replication: Evolution Lifts It, Always}. Le théorème \tierA{} :
en sélection haploïde à fitness constantes, sans biais de transmission, la fitness moyenne ne décroît pas,
$\bar w\le\bar w'$ ; c'est un cas particulier du théorème fondamental de Fisher, déduit de l'équation de Price.
Images \tierX{}. (a)~Trois types, $w=(1;\,1{,}45;\,1{,}9)$ : sur le simplexe coloré par $\bar w$, six trajectoires
de $x_i'=x_iw_i/\bar w$ montent vers le meilleur type. (b)~Cas diploïde sous Hardy--Weinberg, avec surdominance
($w_{11}=0{,}78$, $w_{12}=1$, $w_{22}=0{,}62$) : les deux trajectoires montent sur $\bar w(p)$ jusqu'à
$p^*=0{,}6333$. Ce panneau relève du théorème diploïde \lean{wbar_increment} du chapitre~4 (H12) \tierA{}, pas de
\lean{AURELIA}. (c)~$\bar w$ au fil des générations, jamais décroissante. L'identité
$\bar w'-\bar w=\mathrm{Var}(w)/\bar w$ est vérifiée en rationnels exacts sur 400 cas aléatoires ; contrôle
négatif : avec un biais de mutation du meilleur type vers le moins bon, $\bar w$ baisse de $0{,}35$.}
\label{fig:egg:AURELIA}
\end{figure}
).
% Figures produites par code/easter_egg_images.py (matplotlib/numpy/scipy) ; chaque grandeur tracée est calculée
% à partir de sa formule. Les auto-contrôles numériques (contrôles positifs et négatifs) sont consignés dans
% results/textbook_formalisation_2026/easter_eggs_selfcheck.json ; un contrôle en échec interrompt le script.
\section*{Les cinq images}
\addcontentsline{toc}{section}{Les cinq images}

Chacun des cinq noms courts de la section~\ref{ch5:acronymes} reçoit ici une image. Le niveau du théorème
(\tierA{}) porte sur l'énoncé Lean ré-exporté. Les courbes, les couleurs et les valeurs affichées sont des calculs
numériques \tierX{} : elles illustrent un cas particulier de l'énoncé, sans le démontrer ni l'étendre.

\begin{figure}[htbp]
\centering
\includegraphics[width=\textwidth]{egg_ALIX.pdf}
\caption{\lean{ALIX}, \emph{A Lattice Is eXtremal}. Le théorème \tierA{} (conditionnel au théorème amont
\lean{universal_energy_minimum} d'openai/math, généré par IA et non relu par des pairs) dit qu'à toute densité
$\rho>0$ le réseau triangulaire minimise l'énergie de tout potentiel admissible, parmi toutes les configurations
localement finies de densité~$\rho$. L'image \tierX{} prend le noyau gaussien $e^{-\pi\alpha r^2}$, $\alpha=0{,}4$,
à densité~1, et ne compare que des \emph{réseaux}. (a)~et~(b) : cellules de Voronoï colorées par la contribution
$e^{-\pi\alpha|p|^2}$ de chaque site~$p$ à l'énergie du site central ; $E_\alpha=1{,}50173$ (triangulaire)
contre $1{,}50388$ (carré). (c)~: écart $E_\alpha(\tau)-E_\alpha(e^{i\pi/3})$ sur le domaine fondamental des formes
de réseau, calculé par sommes thêta ; le minimum de la grille (au point équivalent $-\tfrac12+i\tfrac{\sqrt3}2$, à un pas de grille près) et cinq minimisations de Nelder--Mead (à $10^{-7}$ près) retombent sur
$\tau=e^{i\pi/3}$. En médaillon : $E_{\text{carré}}/E_{\text{tri}}-1>0$ pour
$\alpha\in[0{,}15;6]$, c'est-à-dire pour toute une gamme de densités.}
\label{fig:egg:ALIX}
\end{figure}

\begin{figure}[htbp]
\centering
\includegraphics[width=\textwidth]{egg_KAL.pdf}
\caption{\lean{KAL}, \emph{Kolmogorov Absorption Limit}. Le théorème \tierA{} : pour une chaîne sur
$\{0,\dots,N\}$ à bords absorbants, si $h$ est harmonique aux états intérieurs, $h(0)=0$, $h(N)=1$, et si la
variance à un pas de $h$ est strictement positive, alors $P^n(i,N)\to h(i)$ et $P^n(i,0)\to 1-h(i)$. Images
\tierX{}, par puissances exactes de la matrice de transition. (a)~Wright--Fisher neutre, $N=24$ : $P^n(i,N)$
converge vers $h(i)=i/N$ (bande de droite) ; écart maximal à $n=4096$ : $6\cdot10^{-16}$. (b)~Moran avec
sélection, $N=24$, $r=1{,}25$ : $P^n(i,N)$ tend vers $(1-r^{-i})/(1-r^{-N})$ (pointillés), écart
$7\cdot10^{-15}$ à $n=2^{22}$. L'harmonicité des deux $h$ est vérifiée en rationnels exacts (pour Wright--Fisher, à $N=8$).}
\label{fig:egg:KAL}
\end{figure}

\begin{figure}[htbp]
\centering
\includegraphics[width=\textwidth]{egg_ILENE.pdf}
\caption{\lean{ILENE}, \emph{Isothermal Law: Every Node Equivalent}. Le théorème \tierA{} : sur un graphe
isotherme ($W$ bistochastique) fortement connexe, en naissance--mort à sélection uniforme, un mutant unique de
fitness $r\neq1$ se fixe avec la probabilité $(1-r^{-1})/(1-r^{-n})$, quel que soit son sommet de départ. Images
\tierX{}, $r=2$, probabilités obtenues en résolvant le système linéaire sur les $2^n$ configurations (en flottants pour $n=8$, en rationnels pour $n=3$), et
contrôlées par la limite de $P^t$. (a)~Graphe pondéré à $n=8$ sommets, $W=\frac12\Pi_1+\frac13\Pi_2+\frac16\Pi_3$
(trois permutations sans point fixe) : les huit sommets donnent la valeur de Moran $0{,}50196$ à $2\cdot10^{-15}$ près.
(b)~L'étoile à 7 feuilles n'est pas isotherme et sort des hypothèses : centre $0{,}167$, feuilles $0{,}700$.
(c)~Calcul en rationnels pour $n=3$ : le 3-cycle donne $\frac47$ partout ; l'étoile donne $\frac5{12}$ au centre et
$\frac23$ aux feuilles, de moyenne $\frac7{12}\neq\frac47$.}
\label{fig:egg:ILENE}
\end{figure}

\begin{figure}[htbp]
\centering
\includegraphics[width=\textwidth]{egg_ALEXIS.pdf}
\caption{\lean{ALEXIS}, \emph{Averaged Laplace EXponentials Inherit Signs}. Le théorème \tierA{} : un mélange
positif d'exponentielles $t\mapsto\int e^{-t\,k(x)}m(x)\,d\mu(x)$ est un potentiel admissible, de classe
$C^\infty$ sur $]0,+\infty[$, positif, avec $(-1)^n g^{(n)}\ge0$. L'image \tierX{} en montre une instance, le
potentiel de Rytova--Keldysh ($r_0=1$), écrit comme mélange de gaussiennes
$V(r)=\int_0^\infty \frac{w(a)}{2a}\,e^{-r^2/4a}\,da$ de poids $w(a)=\frac1{r_0}e^{a/r_0^2}\operatorname{erfc}
(\sqrt a/r_0)>0$. (a)~Les contributions de douze tranches en $a$, empilées, redonnent la forme de Struve
$\frac{\pi}{2r_0}[\mathbf H_0-Y_0](r/r_0)$ (points) à moins de $3\cdot10^{-13}$ près en relatif ; l'identification avec la
forme de Struve est une preuve humaine, non vérifiée par machine \tierC{}. (b)~Le poids $w$, coloré par tranche.
(c)~$(-1)^n g^{(n)}(t)$, $g(t)=V(\sqrt t)$, pour $n=0,\dots,6$ : toutes les courbes sont positives, donc les
dérivées $g^{(n)}$ alternent de signe ; $n=1$ est recoupé par différences finies de la forme de Struve.}
\label{fig:egg:ALEXIS}
\end{figure}

\begin{figure}[htbp]
\centering
\includegraphics[width=\textwidth]{egg_AURELIA.pdf}
\caption{\lean{AURELIA}, \emph{Average Under Replication: Evolution Lifts It, Always}. Le théorème \tierA{} :
en sélection haploïde à fitness constantes, sans biais de transmission, la fitness moyenne ne décroît pas,
$\bar w\le\bar w'$ ; c'est un cas particulier du théorème fondamental de Fisher, déduit de l'équation de Price.
Images \tierX{}. (a)~Trois types, $w=(1;\,1{,}45;\,1{,}9)$ : sur le simplexe coloré par $\bar w$, six trajectoires
de $x_i'=x_iw_i/\bar w$ montent vers le meilleur type. (b)~Cas diploïde sous Hardy--Weinberg, avec surdominance
($w_{11}=0{,}78$, $w_{12}=1$, $w_{22}=0{,}62$) : les deux trajectoires montent sur $\bar w(p)$ jusqu'à
$p^*=0{,}6333$. Ce panneau relève du théorème diploïde \lean{wbar_increment} du chapitre~4 (H12) \tierA{}, pas de
\lean{AURELIA}. (c)~$\bar w$ au fil des générations, jamais décroissante. L'identité
$\bar w'-\bar w=\mathrm{Var}(w)/\bar w$ est vérifiée en rationnels exacts sur 400 cas aléatoires ; contrôle
négatif : avec un biais de mutation du meilleur type vers le moins bon, $\bar w$ baisse de $0{,}35$.}
\label{fig:egg:AURELIA}
\end{figure}
).
% Figures produites par code/easter_egg_images.py (matplotlib/numpy/scipy) ; chaque grandeur tracée est calculée
% à partir de sa formule. Les auto-contrôles numériques (contrôles positifs et négatifs) sont consignés dans
% results/textbook_formalisation_2026/easter_eggs_selfcheck.json ; un contrôle en échec interrompt le script.
\section*{Les cinq images}
\addcontentsline{toc}{section}{Les cinq images}

Chacun des cinq noms courts de la section~\ref{ch5:acronymes} reçoit ici une image. Le niveau du théorème
(\tierA{}) porte sur l'énoncé Lean ré-exporté. Les courbes, les couleurs et les valeurs affichées sont des calculs
numériques \tierX{} : elles illustrent un cas particulier de l'énoncé, sans le démontrer ni l'étendre.

\begin{figure}[htbp]
\centering
\includegraphics[width=\textwidth]{egg_ALIX.pdf}
\caption{\lean{ALIX}, \emph{A Lattice Is eXtremal}. Le théorème \tierA{} (conditionnel au théorème amont
\lean{universal_energy_minimum} d'openai/math, généré par IA et non relu par des pairs) dit qu'à toute densité
$\rho>0$ le réseau triangulaire minimise l'énergie de tout potentiel admissible, parmi toutes les configurations
localement finies de densité~$\rho$. L'image \tierX{} prend le noyau gaussien $e^{-\pi\alpha r^2}$, $\alpha=0{,}4$,
à densité~1, et ne compare que des \emph{réseaux}. (a)~et~(b) : cellules de Voronoï colorées par la contribution
$e^{-\pi\alpha|p|^2}$ de chaque site~$p$ à l'énergie du site central ; $E_\alpha=1{,}50173$ (triangulaire)
contre $1{,}50388$ (carré). (c)~: écart $E_\alpha(\tau)-E_\alpha(e^{i\pi/3})$ sur le domaine fondamental des formes
de réseau, calculé par sommes thêta ; le minimum de la grille (au point équivalent $-\tfrac12+i\tfrac{\sqrt3}2$, à un pas de grille près) et cinq minimisations de Nelder--Mead (à $10^{-7}$ près) retombent sur
$\tau=e^{i\pi/3}$. En médaillon : $E_{\text{carré}}/E_{\text{tri}}-1>0$ pour
$\alpha\in[0{,}15;6]$, c'est-à-dire pour toute une gamme de densités.}
\label{fig:egg:ALIX}
\end{figure}

\begin{figure}[htbp]
\centering
\includegraphics[width=\textwidth]{egg_KAL.pdf}
\caption{\lean{KAL}, \emph{Kolmogorov Absorption Limit}. Le théorème \tierA{} : pour une chaîne sur
$\{0,\dots,N\}$ à bords absorbants, si $h$ est harmonique aux états intérieurs, $h(0)=0$, $h(N)=1$, et si la
variance à un pas de $h$ est strictement positive, alors $P^n(i,N)\to h(i)$ et $P^n(i,0)\to 1-h(i)$. Images
\tierX{}, par puissances exactes de la matrice de transition. (a)~Wright--Fisher neutre, $N=24$ : $P^n(i,N)$
converge vers $h(i)=i/N$ (bande de droite) ; écart maximal à $n=4096$ : $6\cdot10^{-16}$. (b)~Moran avec
sélection, $N=24$, $r=1{,}25$ : $P^n(i,N)$ tend vers $(1-r^{-i})/(1-r^{-N})$ (pointillés), écart
$7\cdot10^{-15}$ à $n=2^{22}$. L'harmonicité des deux $h$ est vérifiée en rationnels exacts (pour Wright--Fisher, à $N=8$).}
\label{fig:egg:KAL}
\end{figure}

\begin{figure}[htbp]
\centering
\includegraphics[width=\textwidth]{egg_ILENE.pdf}
\caption{\lean{ILENE}, \emph{Isothermal Law: Every Node Equivalent}. Le théorème \tierA{} : sur un graphe
isotherme ($W$ bistochastique) fortement connexe, en naissance--mort à sélection uniforme, un mutant unique de
fitness $r\neq1$ se fixe avec la probabilité $(1-r^{-1})/(1-r^{-n})$, quel que soit son sommet de départ. Images
\tierX{}, $r=2$, probabilités obtenues en résolvant le système linéaire sur les $2^n$ configurations (en flottants pour $n=8$, en rationnels pour $n=3$), et
contrôlées par la limite de $P^t$. (a)~Graphe pondéré à $n=8$ sommets, $W=\frac12\Pi_1+\frac13\Pi_2+\frac16\Pi_3$
(trois permutations sans point fixe) : les huit sommets donnent la valeur de Moran $0{,}50196$ à $2\cdot10^{-15}$ près.
(b)~L'étoile à 7 feuilles n'est pas isotherme et sort des hypothèses : centre $0{,}167$, feuilles $0{,}700$.
(c)~Calcul en rationnels pour $n=3$ : le 3-cycle donne $\frac47$ partout ; l'étoile donne $\frac5{12}$ au centre et
$\frac23$ aux feuilles, de moyenne $\frac7{12}\neq\frac47$.}
\label{fig:egg:ILENE}
\end{figure}

\begin{figure}[htbp]
\centering
\includegraphics[width=\textwidth]{egg_ALEXIS.pdf}
\caption{\lean{ALEXIS}, \emph{Averaged Laplace EXponentials Inherit Signs}. Le théorème \tierA{} : un mélange
positif d'exponentielles $t\mapsto\int e^{-t\,k(x)}m(x)\,d\mu(x)$ est un potentiel admissible, de classe
$C^\infty$ sur $]0,+\infty[$, positif, avec $(-1)^n g^{(n)}\ge0$. L'image \tierX{} en montre une instance, le
potentiel de Rytova--Keldysh ($r_0=1$), écrit comme mélange de gaussiennes
$V(r)=\int_0^\infty \frac{w(a)}{2a}\,e^{-r^2/4a}\,da$ de poids $w(a)=\frac1{r_0}e^{a/r_0^2}\operatorname{erfc}
(\sqrt a/r_0)>0$. (a)~Les contributions de douze tranches en $a$, empilées, redonnent la forme de Struve
$\frac{\pi}{2r_0}[\mathbf H_0-Y_0](r/r_0)$ (points) à moins de $3\cdot10^{-13}$ près en relatif ; l'identification avec la
forme de Struve est une preuve humaine, non vérifiée par machine \tierC{}. (b)~Le poids $w$, coloré par tranche.
(c)~$(-1)^n g^{(n)}(t)$, $g(t)=V(\sqrt t)$, pour $n=0,\dots,6$ : toutes les courbes sont positives, donc les
dérivées $g^{(n)}$ alternent de signe ; $n=1$ est recoupé par différences finies de la forme de Struve.}
\label{fig:egg:ALEXIS}
\end{figure}

\begin{figure}[htbp]
\centering
\includegraphics[width=\textwidth]{egg_AURELIA.pdf}
\caption{\lean{AURELIA}, \emph{Average Under Replication: Evolution Lifts It, Always}. Le théorème \tierA{} :
en sélection haploïde à fitness constantes, sans biais de transmission, la fitness moyenne ne décroît pas,
$\bar w\le\bar w'$ ; c'est un cas particulier du théorème fondamental de Fisher, déduit de l'équation de Price.
Images \tierX{}. (a)~Trois types, $w=(1;\,1{,}45;\,1{,}9)$ : sur le simplexe coloré par $\bar w$, six trajectoires
de $x_i'=x_iw_i/\bar w$ montent vers le meilleur type. (b)~Cas diploïde sous Hardy--Weinberg, avec surdominance
($w_{11}=0{,}78$, $w_{12}=1$, $w_{22}=0{,}62$) : les deux trajectoires montent sur $\bar w(p)$ jusqu'à
$p^*=0{,}6333$. Ce panneau relève du théorème diploïde \lean{wbar_increment} du chapitre~4 (H12) \tierA{}, pas de
\lean{AURELIA}. (c)~$\bar w$ au fil des générations, jamais décroissante. L'identité
$\bar w'-\bar w=\mathrm{Var}(w)/\bar w$ est vérifiée en rationnels exacts sur 400 cas aléatoires ; contrôle
négatif : avec un biais de mutation du meilleur type vers le moins bon, $\bar w$ baisse de $0{,}35$.}
\label{fig:egg:AURELIA}
\end{figure}
